
%(BEGIN_QUESTION)
% Copyright 2006, Tony R. Kuphaldt, released under the Creative Commons Attribution License (v 1.0)
% This means you may do almost anything with this work of mine, so long as you give me proper credit

Forklar hvorfor glassmåleelektroden i en pH-sensor alltid må holdes fuktig, ideelt i kaliumklorid-løsning (KCl), under lagring.

\underbar{file i00620}
%(END_QUESTION)





%(BEGIN_ANSWER)

Funksjonen til elektroden avhenger av at glasset har indre og ytre {\it hydrerte} lag. Hvis elektroden tørker på utsiden, kan ioner ikke permeere, og elektroden blir ubrukelig.

De hydrerte lagene løses langsomt opp over tid, slik at nye (tørre) glasslag under hydreres og overtar funksjonen. Vedlikehold av dette hydrerte laget er kritisk for korrekt drift.

Dette forklarer også hvorfor pH-prober har begrenset levetid, og hvorfor levetiden reduseres ved høy temperatur. Når temperaturen øker, øker oppløsningshastigheten i glasset, og elektroden slites raskere.

%(END_ANSWER)





%(BEGIN_NOTES)


%INDEX% Måling, analytisk: pH

%(END_NOTES)

